\chapter{Commentator And Translator Comparison Table}

\section{Purpose}

This appendix summarizes the main comparison points used by the argument. It is not a substitute for page-level citation. Its role is to show the kinds of interpretive decisions that translations and commentators make in Revelation 1:1--3.

\section{Translation Pressure Points}

\begin{center}
\begin{longtable}{@{}L{0.2\textwidth}L{0.34\textwidth}L{0.34\textwidth}@{}}
\toprule
Pressure point & Common smooth result & Question reopened by this book \\
\midrule
\textit{apokalypsis} & ``Revelation'' or title-like label & Is the term naming the written book, an eventive disclosure, or both in different relations? \\
\textit{Iesou Christou} & from Jesus Christ, about Jesus Christ, or broad ``of Jesus Christ'' & What relation does the structure validate rather than merely permit? \\
\textit{auto} & Jesus Christ as recipient & Does the B / B' transfer pair identify John instead? \\
\textit{deixai} & Jesus Christ shows the servants & Is the subject inherited from the linear chain or governed by the larger structure? \\
\textit{tois doulois autou} & Christians generally & How does servant language function inside the disclosure system? \\
\textit{en tachei} & soon or shortly & Does the phrase function as a timetable, a speed/manner phrase, or a more complex timing claim? \\
\textit{esemanen} & made known & Does signifying remain distinct from showing and ordinary informing? \\
\textit{aposteilas} & by sending his angel to John & What is sent, who sends, and how does sending correspond to giving? \\
\textit{hosa eiden} & everything John saw & How does sight relate to the word of God and testimony of Jesus Christ? \\
\bottomrule
\end{longtable}
\end{center}

\section{Representative Commentary Patterns}

\begin{center}
\begin{longtable}{@{}L{0.24\textwidth}L{0.31\textwidth}L{0.35\textwidth}@{}}
\toprule
Pattern & What it usually supplies & What the present argument tests \\
\midrule
Linear chain & God gives to Jesus Christ; Jesus Christ communicates through an angel to John; John writes for the churches. & Whether each step is required by the surface form or produced by inherited smoothness. \\
Hidden-revelation emphasis & \textit{Apokalypsis} is treated as disclosure of previously hidden divine mysteries or future events. & Whether hiddenness should govern before the A / A' testimony relation is considered. \\
Broad servant audience & The servants are identified with believers or Christian readers generally. & Whether the servant category is governed by reception, obedience, and John's servant role. \\
Temporal imminence & \textit{En tachei} is taken mainly as soon or shortly. & Whether speed and temporal nearness need to be distinguished before the phrase controls the whole book. \\
Signifying flattened & \textit{Esemanen} is rendered as made known or communicated. & Whether signifying names a marked mode of communication paired with showing. \\
Vision-summary reading & John testifies to the word/testimony, understood as what he saw. & Whether the final visual clause should summarize, supplement, or be integrated with the named testimony objects. \\
\bottomrule
\end{longtable}
\end{center}

\section{Citation Status}

The comparison above is a map of interpretive patterns. Named sources should be cited in the main chapters only when their exact page or stable online location has been verified in the page-reference ledger. Until that verification is complete, the main text should continue to use general phrases such as ``many translations,'' ``a common commentary reading,'' and ``the familiar reading.''
